\documentclass[conference,a4paper]{IEEEtran}
\hyphenpenalty=5000
\tolerance=1000
\usepackage[utf8]{inputenc}
\usepackage[T1]{fontenc}
\usepackage[spanish,english]{babel}
\usepackage{graphicx}
\usepackage{textcomp} % For special characters

% We'll use direct Unicode instead of \DeclareUnicodeCharacter
% since we're using lualatex which supports Unicode directly

% Información del artículo
\title{Evaluación de algoritmos bioinspirados recientes para el Vehicle Routing Problem
\\ \textit{Evaluation of Recent Bio-inspired Algorithms for the Vehicle Routing Problem}}

\author{%
  \begin{tabular}{c}
    \textbf{Felipe Ignacio González G.}\\
    Programa de Magíster en Informática\\
    Universidad de Valparaíso, Chile\\
    felipe.gonzalezggo@postgrado.uv.cl
  \end{tabular}}

\begin{document}
\selectlanguage{spanish}
\maketitle

% Resumen en español
\begin{abstract}
Se presentó un estudio comparativo riguroso de algoritmos metaheurísticos bioinspirados aplicados al Problema de Ruteo de Vehículos (VRP).
\end{abstract}

\begin{IEEEkeywords}
metaheurísticas bioinspiradas, vehicle routing problem
\end{IEEEkeywords}

% Abstract en inglés
\selectlanguage{english}
\begin{abstract}
This paper presents a rigorous comparative study of bio-inspired metaheuristic algorithms applied to the Vehicle Routing Problem (VRP).
\end{abstract}

\begin{IEEEkeywords}
bio-inspired metaheuristics, vehicle routing problem
\end{IEEEkeywords}

\selectlanguage{spanish}

% Introducción
\section{Introducción}
El Problema de Ruteo de Vehículos (VRP) constituye uno de los desafíos fundamentales en investigación operativa y optimización combinatoria.

\section{Conclusiones}
Este estudio presentó una evaluación comparativa rigurosa de algoritmos metaheurísticos bioinspirados aplicados al problema VRP.

\end{document}